\problema{Number of Provinces}{547}{src/01-grafos/547-number-of-provinces.cpp}{M}

Hay \texttt{n} ciudades y una matriz de adyacencia \texttt{isConnected} de
tamaño \texttt{n x n}, donde \texttt{isConnected[i][j] == 1} si la ciudad
\texttt{i} y la ciudad \texttt{j} están conectadas \textbf{directamente}
(la matriz es simétrica y siempre \texttt{isConnected[i][i] == 1}). Una
\emph{provincia} es un grupo de ciudades conectadas directa o
indirectamente, sin conexión con ninguna ciudad fuera del grupo. Queremos
contar cuántas provincias hay.

\paragraph{Intuición.} Esta matriz no es más que otra forma de entregarnos
el mismo grafo no dirigido de \textbf{323. Number of Connected
Components}: en vez de recibir una lista de aristas $[a, b]$, nos dan,
para cada par de ciudades $(i, j)$, si existe o no esa arista. El
problema de fondo es idéntico: contar componentes conexos. Basta
recorrer la matriz y, cada vez que \texttt{isConnected[i][j] == 1} con
\texttt{i != j}, tratar a $i$ y $j$ como dos nodos unidos por una arista.

\begin{center}
\ttfamily
\begin{tabular}{ccc}
1 & 1 & 0\\
1 & 1 & 0\\
0 & 0 & 1\\
\end{tabular}
\end{center}

\noindent En esta matriz hay \textbf{2} provincias: $\{0, 1\}$ (conectadas
directamente) y $\{2\}$ (aislada, solo conectada consigo misma en la
diagonal).

\paragraph{Solución (Union-Find / DSU).} Igual que en 323, usamos
Union-Find: cada ciudad empieza siendo su propia raíz (\texttt{n}
provincias). Recorremos únicamente la mitad superior de la matriz (es
simétrica, así que revisar $i < j$ basta) y, cuando
\texttt{isConnected[i][j] == 1}, unimos las raíces de $i$ y $j$ con
compresión de caminos y unión por rango; si las raíces eran distintas,
dos provincias acaban de fusionarse y decrementamos el contador.

\lstinputlisting{src/01-grafos/547-number-of-provinces.cpp}

\complejidad{$O(n^2\cdot\alpha(n))$}{$O(n)$}

\noindent recorremos las $n^2$ celdas de la matriz (en la práctica, solo
la mitad superior), y cada \texttt{find}/\texttt{unir} es casi $O(1)$
amortizado gracias a $\alpha(n)$, la inversa de Ackermann.

\paragraph{Notas de entrevista (Amazon).} Este problema es hermano
directo de \textbf{323. Number of Connected Components}: la única
diferencia real es el formato de entrada (matriz de adyacencia en vez de
lista de aristas), la técnica de fondo es la misma Union-Find. También se
puede resolver con DFS/BFS clásico: por cada ciudad no visitada, sumar
una provincia y recorrer con DFS toda la fila de la matriz para marcar
como visitadas a todas las ciudades conectadas (transitivamente), igual
que se ``hunde'' la tierra en \textbf{200. Number of Islands}. Vale la
pena mencionar ambos enfoques: DFS/BFS es más directo de explicar dado
que la entrada ya es una matriz densa, pero Union-Find generaliza mejor
si el entrevistador hace un seguimiento tipo ``¿y si las conexiones
llegan una a una en vez de venir todas de entrada?''. Cuidado con el
costo: si $n$ es grande, $O(n^2)$ ya es el costo de solo \emph{leer} la
matriz de adyacencia, así que ninguna técnica baja de esa cota; con una
lista de adyacencia dispersa convendría más recibir aristas explícitas
en vez de la matriz completa. Casos borde: una sola ciudad ($1$
provincia), ninguna ciudad conectada con otra (cada una su propia
provincia), y todas las ciudades conectadas entre sí ($1$ provincia).
